%! Exponential Functions — Unit 4, Lesson 4.3 — Guided Notes
\documentclass[10pt]{article}
\usepackage{precalculus-article}
\usepackage{precalculus-boxes}
\newcommand{\TallMath}[1]{$\displaystyle #1\rule[-1.4em]{0pt}{3.2em}$}

\begin{document}

\pageheader{Unit 4, Lesson 4.3}{Guided Notes}
\namedateperiod

\vspace{0.06in}

\begin{objectivebox}
By the end of this lesson, I will be able to\ldots
\begin{enumerate}[itemsep=2pt]
  \item Identify the initial value $a$ and base $b$ of $f(x) = ab^x$ and classify as
        \blank{3cm} growth or \blank{3cm} decay. \hfill\textit{(LO 2.3.A.1)}
  \item Describe the \blank{2.5cm}, end \blank{2.5cm}, \blank{3cm}, and monotonic
        nature of an exponential function. \hfill\textit{(LO 2.3.A.3 / A.5)}
  \item Determine whether a function is exponential by checking whether outputs are
        \blank{4cm} over equal-length input intervals. \hfill\textit{(LO 2.3.A.3 / A.4)}
\end{enumerate}
\end{objectivebox}

\vspace{0.04in}

\begin{vocabbox}
\termblanklong{Exponential function}
\termblanklong{Initial value}
\termblanklong{Base}
\termblanklong{Exponential growth}
\termblanklong{Exponential decay}
\termblanklong{Concave up / concave down}
\termblanklong{End behavior}
\termblanklong{Horizontal asymptote}
\end{vocabbox}

\vspace{0.04in}

\begin{hookbox}
\textbf{The Doubling Penny:} If you receive 1 penny on Day 1, 2 on Day 2, 4 on Day 3, \ldots,
how much total after 30 days?

\medskip
My guess: \blank{4cm} \quad Actual: $2^{30} - 1 \approx \$\,10.7$ million.

\medskip
The function modeling Day $x$ is $f(x) = 1 \cdot \blank{1cm}^x$.
The initial value is \blank{1.5cm} and the base is \blank{1.5cm}.
\end{hookbox}

\vspace{0.04in}

\begin{notesbox}{LO 2.3.A.1 --- General Form of an Exponential Function (EK 2.3.A.1)}
\[
  f(x) = a \cdot b^x, \quad a \ne 0,\ b > 0,\ b \ne 1
\]

\begin{itemize}[itemsep=3pt]
  \item $a = f(0) = $ \blank{3cm} (the $y$-intercept / initial value)
  \item $b = $ \blank{5cm} (the base; applied once per unit input increase)
  \item \textbf{Exponential growth:} $a > 0$ and $b$ \blank{1.5cm} 1
  \item \textbf{Exponential decay:} $a > 0$ and \blank{2cm} $< b <$ \blank{1cm}
\end{itemize}

\medskip
\textbf{Example:} $f(x) = 3 \cdot 2^x$. \quad $a = $ \blank{1.5cm} \quad $b = $ \blank{1.5cm}
\quad Growth or decay? \blank{3cm}

\medskip
\renewcommand{\arraystretch}{1.5}
\begin{tabular}{|c|c|c|c|c|}
\hline
$x$    & 0 & 1 & 2 & 3 \\
\hline
$f(x)$ & \blank{1.2cm} & \blank{1.2cm} & \blank{1.2cm} & \blank{1.2cm} \\
\hline
\multicolumn{5}{|l|}{Ratio of successive outputs: \quad $\dfrac{f(1)}{f(0)} = $ \blank{1.5cm}
\quad $\dfrac{f(2)}{f(1)} = $ \blank{1.5cm} \quad $\dfrac{f(3)}{f(2)} = $ \blank{1.5cm}} \\
\hline
\end{tabular}
\end{notesbox}

\vspace{0.04in}

\begin{notesbox}{LO 2.3.A.2/A.3 --- Key Characteristics (EK 2.3.A.2, A.3)}
\renewcommand{\arraystretch}{1.6}
\begin{tabularx}{\linewidth}{@{} >{\bfseries}p{0.28\linewidth} X @{}}
Domain & \blank{7cm} \\
Range (when $a > 0$) & \blank{5cm}; horizontal asymptote: $y = $ \blank{1.5cm} \\
Monotonicity & Always \blank{3cm} if $b > 1$; always \blank{3cm} if $0 < b < 1$ \\
Extrema on open intervals & \blank{6cm} \\
Concavity (when $a > 0$) & Always concave \blank{2.5cm}; no \blank{3cm} points \\
\end{tabularx}
\end{notesbox}

\vspace{0.04in}

\begin{notesbox}{LO 2.3.A.5 --- End Behavior (EK 2.3.A.5)}
For $f(x) = ab^x$ with $a > 0$, $b > 1$ (growth):
\[
  \lim_{x \to \infty} f(x) = \blank{2cm}
  \qquad
  \lim_{x \to -\infty} f(x) = \blank{2cm}
\]

For $f(x) = ab^x$ with $a > 0$, $0 < b < 1$ (decay):
\[
  \lim_{x \to \infty} f(x) = \blank{2cm}
  \qquad
  \lim_{x \to -\infty} f(x) = \blank{2cm}
\]

\textbf{Example:} State the end behavior of $h(x) = 5 \cdot (0.4)^x$.

\vspace{0.04in}
$\displaystyle\lim_{x\to\infty} h(x) = $ \blank{2cm} \qquad
$\displaystyle\lim_{x\to-\infty} h(x) = $ \blank{2cm}
\end{notesbox}

\vspace{0.04in}

\begin{notesbox}{LO 2.3.A.4 --- Detecting Exponential from a Transformed Table (EK 2.3.A.4)}
\textbf{Key idea:} If $g(x) = f(x) + k$ and $f$ is exponential, the outputs of $g$ will
\textit{not} be proportional. To detect $f$: subtract \blank{1.5cm} from each output of $g$,
then check whether the resulting values have a constant \blank{2cm}.

\medskip
\textbf{Example:} Outputs of $g(x)$: 7, 10, 16, 28 (inputs 0, 1, 2, 3). Suppose $k = 4$.

\vspace{0.04in}
Subtract 4: \blank{1.2cm}, \blank{1.2cm}, \blank{1.2cm}, \blank{1.2cm}

Ratios: \blank{1.5cm}, \blank{1.5cm}, \blank{1.5cm}

Are the ratios constant? \blank{2cm} \quad Write $f(x) = $ \blank{3cm}
\quad Write $g(x) = $ \blank{4cm}
\end{notesbox}

\end{document}
